% Uniaxial in-plane stretching of a rectangular plate — analytical derivation
% AIM Interactive Mechanics series. Compile with pdflatex.
\documentclass[10pt,a4paper]{article}
\usepackage[margin=2cm]{geometry}
\usepackage{amsmath}
\usepackage[T1]{fontenc}
\usepackage{lmodern}
\usepackage{microtype}
\usepackage{titlesec}
\titlespacing*{\section}{0pt}{9pt}{2pt}
\usepackage{xcolor}
\definecolor{iron}{HTML}{935636}
\definecolor{celdark}{HTML}{4B6C61}
\definecolor{celline}{HTML}{CDE0D8}
\usepackage[colorlinks=true,urlcolor=iron,linkcolor=black]{hyperref}
\setlength{\parindent}{0pt}
\setlength{\parskip}{4pt}
\AtBeginDocument{\setlength{\abovedisplayskip}{6pt}\setlength{\belowdisplayskip}{6pt}\setlength{\abovedisplayshortskip}{3pt}\setlength{\belowdisplayshortskip}{3pt}}
\pagestyle{empty}

\begin{document}


{\LARGE Uniaxial Plate Stretching}

\medskip
{\color{celline}\hrule height 1pt}

\section*{Setup and assumptions}
A thin rectangular plate of length $a$, width $b$ and thickness $t \ll a, b$ is
loaded in its plane by a uniform tensile traction on the two edges
$x = \pm a/2$, of resultant $F$, so the applied stress is
$\sigma_0 = F/(b\,t)$. The side edges $y = \pm b/2$ and both faces are
traction-free. The material is linear elastic and isotropic ($E$, $\nu$);
strains are small.

\section*{Plane stress}
Because the plate is thin and its faces are free, the through-thickness stresses
cannot build up: $\sigma_{zz} = \tau_{xz} = \tau_{yz} \approx 0$. Hooke's law
then reduces to the plane-stress relations
\begin{equation}
\varepsilon_x = \frac{\sigma_x - \nu \sigma_y}{E}, \qquad
\varepsilon_y = \frac{\sigma_y - \nu \sigma_x}{E}, \qquad
\gamma_{xy} = \frac{\tau_{xy}}{G}, \qquad
\varepsilon_z = -\frac{\nu\,(\sigma_x + \sigma_y)}{E}.
\end{equation}
Note $\varepsilon_z \neq 0$: the plate is free to thin. (Plane \emph{strain},
where $\varepsilon_z = 0$ is enforced, is a different idealisation.)

\section*{The homogeneous field solves the problem exactly}
Try the uniform state
\begin{equation}
\sigma_x = \sigma_0, \qquad \sigma_y = 0, \qquad \tau_{xy} = 0
\qquad \text{everywhere.}
\end{equation}
Check it against everything the exact 2D theory requires:
\begin{itemize}
\setlength\itemsep{2pt}
\item \emph{Equilibrium} $\partial\sigma_x/\partial x + \partial\tau_{xy}/\partial y = 0$
      and its companion: satisfied trivially, all fields constant.
\item \emph{Boundary conditions}: on $x = \pm a/2$ the traction is
      $\sigma_x = \sigma_0$ as applied; on $y = \pm b/2$ the traction components
      $\sigma_y$ and $\tau_{xy}$ vanish as required. Satisfied \emph{pointwise},
      not just on average.
\item \emph{Compatibility}: the strains are constant, so a displacement field
      exists; explicitly, measuring from the plate centre,
      $u = \varepsilon_x\, x$ and $v = \varepsilon_y\, y$.
\end{itemize}
By uniqueness of linear elastic solutions, this \emph{is} the solution --- no
approximation involved.

\section*{Results}
\begin{equation}
\boxed{\;
\varepsilon_x = \frac{\sigma_0}{E}, \qquad
\varepsilon_y = -\nu\,\varepsilon_x, \qquad
\varepsilon_z = -\nu\,\varepsilon_x
\;}
\qquad\Rightarrow\qquad
\Delta a = \varepsilon_x\, a, \qquad
\Delta b = \varepsilon_y\, b, \qquad
\Delta t = \varepsilon_z\, t .
\end{equation}
The plate stretches along the load and contracts both across it and through its
thickness; Poisson's ratio is the whole story of the lateral response. At
$\nu = 0$ (cork-like) there is no contraction; at $\nu = 1/2$ (rubber,
incompressible) the volume change
$\varepsilon_x + \varepsilon_y + \varepsilon_z = (1 - 2\nu)\,\varepsilon_x$
vanishes.

\section*{Worked example (simulator defaults)}
$E = 10$~MPa, $\nu = 0.3$, $a = 500$~mm, $b = 250$~mm, $t = 2$~mm, $F = 400$~N:
\begin{equation}
\sigma_0 = \frac{400}{0.25 \times 0.002} = 0.8~\text{MPa}, \qquad
\varepsilon_x = 8\%, \qquad \varepsilon_y = -2.4\%,
\end{equation}
so $\Delta a = 40$~mm, $\Delta b = -6$~mm and $\Delta t = -0.048$~mm. Set the
simulator sliders to these values and compare; then sweep $\nu$ from $0$ to
$0.5$ and watch only the lateral response change.


\end{document}
